% titel: Blatt 7, A2: Kohlenwasserstoffe C2H2 bis C4H4 - Hybridisierung und Bindungen
% skript: Mortimer Kap. 10 (Hybridisierung, Bindungswinkel, Sigma/Pi)

% Teil I: Werkzeuge, kein L\"osungsweg.

Gegenstand sind vier Kohlenwasserstoffe mit kumulierten
C-Doppel-/Dreifachbindungen: Ethin $\mathrm{HCCH}$~($\mathrm{C_2H_2}$),
Ethen $\mathrm{H_2CCH_2}$~($\mathrm{C_2H_4}$), Allen
$\mathrm{H_2CCCH_2}$~($\mathrm{C_3H_4}$) und Butatrien
$\mathrm{H_2CCCCH_2}$~($\mathrm{C_4H_4}$).

\begin{werkzeugkasten}[Hybridisierung und Bindungswinkel]
\begin{itemize}
  \item \textbf{sp$^3$:} 4 Einfachbindungen, Tetraeder, Winkel
        $\approx 109{,}5^\circ$.
  \item \textbf{sp$^2$:} 1 Doppelbindung (1$\sigma$+1$\pi$), trigonal,
        Winkel $\approx 120^\circ$.
  \item \textbf{sp:} 1 Dreifach- oder 2 Doppelbindungen, linear,
        Winkel $180^\circ$.
  \item Pro Bindungstyp zwischen zwei C-Atomen: 1 $\sigma$-Bindung
        (Kopf-überlappung) plus $0$ (Einfach-), $1$ (Doppel-) bzw.\ 2
        ($\pi$-) Bindungen bei Dreifachbindung.
\end{itemize}
\end{werkzeugkasten}

\begin{werkzeugkasten}[sigma-/pi-Gerüst]
Die $\sigma$-Bindungen definieren das räumliche Gerüst (alle Atome in der
``sigma-Ebene'' liegen koplanar). Die $\pi$-Bindungen liegen senkrecht zur
Bindungsachse (p-Orbitale); bei gekreuzten $\pi$-Systemen k\"onnen
H-Atome aus der gemeinsamen Ebene ``herausdrehen''.
\textit{Skript:} Mortimer Kap. 10, ``Hybridisierung''.
\end{werkzeugkasten}

\begin{achtung}[Darauf achten]
\begin{itemize}
  \item Winkel am sp$^2$-C ($\mathrm{CH_2}$) sind $120^\circ$,
        am sp-C $180^\circ$.
  \item Allen ($\mathrm{H_2CCCH_2}$) hat zwei zueinander senkrechte
        $\pi$-Systeme $\rightarrow$ die beiden $\mathrm{CH_2}$-Ebenen sind
        nicht koplanar (alle H nicht in einer Ebene).
  \item Kettenende und Mitte unterscheiden: End-C sp$^2$, zentrale C sp.
\end{itemize}
\end{achtung}
